% GENERATED FILE. Edit canonical lessons, then run npm run generate:books.
% canonical-source-hash: fnv1a64:a1d64329bb86a946
% canonical-lessons: KA-C518-nemisu, KA-C518-vaja-madu, KA-C518-nirudyoga, KA-C518-nivrtti, KA-C518-guttige, KA-R518-review-512-518

\chapter{Hiring and Retiring}
\label{ch:ka-hiring-and-retiring}
\hlchaptermodality{518}

\begin{chapteropening}
\textbf{By the end of this chapter:} \emph{I can say to appoint, to dismiss, unemployment, retirement and a contract.}

Complete the last lesson of chapter 518: I can say to appoint, to dismiss, unemployment, retirement and a contract.
\end{chapteropening}

\section[\texorpdfstring{\kn{ನೇಮಿಸು} (nēmisu)}{ನೇಮಿಸು (nēmisu)}]{\kn{ನೇಮಿಸು} (nēmisu) — to appoint; to hire}

\label{lesson:KA-C518-nemisu}



\begin{quote}
Before the new one: say the Kannada for to achieve, then the Kannada for to gain.
\end{quote}

\subsection*{You'll want to know: \kn{ನೇಮಿಸು}}
\textbf{\kn{ನೇಮಿಸು}} — \emph{nēmisu} — \textquotedblleft{}to appoint; to hire\textquotedblright{}.

\begin{quote}
\kn{ಕಂಪನಿ} \kn{ಹತ್ತು} \kn{ಹೊಸ} \kn{ನೌಕರರನ್ನು} \kn{ನೇಮಿಸಿತು}.
\end{quote}

\begin{quote}
\emph{kampani hattu hosa naukararannu nēmisitu.}
\end{quote}

\begin{quote}
The company hired ten new employees.
\end{quote}

Past (I …): \kn{ನೇಮಿಸಿದೆ}, \emph{nēmiside}.

The one appointed takes -\kn{ಅನ್ನು}. \kn{ನೇಮಕ} is an appointment.

\begin{cousinweb}
\textbf{\kn{ನೇಮಿಸು}} (\emph{nēmisu}) is built on \emph{nēma}, \textquotedblleft{}a fixed rule or arrangement\textquotedblright{}, a Kannada form of Sanskrit \emph{niyama}, \textquotedblleft{}restraint, fixed law\textquotedblright{} (CDIAL 7250), with the verb ending -\emph{isu}.
\end{cousinweb}

\begin{grammarlens}[title={What you've built}]
One more word for describing what you lived and felt.
\end{grammarlens}

\subsection*{Your turn}
\begin{itemize}
\raggedright
  \item \emph{Say it:} \emph{nēmisu}
  \item \emph{Say it:} \emph{nēmisu}, once more
  \item \emph{Say it:} \emph{kampani hattu hosa naukararannu nēmisitu.} — The company hired ten new employees.
  \item \emph{Recall:} say the Kannada for to achieve, then the Kannada for to gain, then say \emph{nēmisu} again
\end{itemize}

\begin{tcolorbox}[breakable,colback=teal!4,colframe=teal!35!black,title={Before you move on}]
What does \kn{ನೇಮಿಸು} mean? (\textquotedblleft{}To appoint; to hire\textquotedblright{}.) Say it once more.
\end{tcolorbox}
\section[\texorpdfstring{\kn{ವಜಾ} \kn{ಮಾಡು} (vajā māḍu)}{ವಜಾ ಮಾಡು (vajā māḍu)}]{\kn{ವಜಾ} \kn{ಮಾಡು} (vajā māḍu) — to dismiss; to sack}

\label{lesson:KA-C518-vaja-madu}



\begin{quote}
Before the new one: say the Kannada for to gain, then the Kannada for to appoint.
\end{quote}

\subsection*{You'll want to know: \kn{ವಜಾ} \kn{ಮಾಡು}}
\textbf{\kn{ವಜಾ} \kn{ಮಾಡು}} — \emph{vajā māḍu} — \textquotedblleft{}to dismiss; to sack\textquotedblright{}.

\begin{quote}
\kn{ತಡವಾಗಿ} \kn{ಬರುತ್ತಿದ್ದ} \kn{ನೌಕರನನ್ನು} \kn{ಕಂಪನಿ} \kn{ವಜಾ} \kn{ಮಾಡಿತು}.
\end{quote}

\begin{quote}
\emph{taḍavāgi baruttidda naukaranannu kampani vajā māḍitu.}
\end{quote}

\begin{quote}
The company sacked the employee who kept coming late.
\end{quote}

Past (I …): \kn{ವಜಾ} \kn{ಮಾಡಿದೆ}, \emph{vajā māḍide}.

Also of a court case: \kn{ಅರ್ಜಿ} \kn{ವಜಾ} \kn{ಆಯಿತು}, the petition was dismissed.

\begin{cousinweb}
\textbf{\kn{ವಜಾ} \kn{ಮಾಡು}} (\emph{vajā māḍu}) joins \emph{vajā}, a word that came into Kannada through the language of offices and courts, with \kn{ಮಾಡು}, \textquotedblleft{}to do\textquotedblright{}; the earlier source of \emph{vajā} is not settled.
\end{cousinweb}

\begin{grammarlens}[title={What you've built}]
One more word for describing what you lived and felt.
\end{grammarlens}

\subsection*{Your turn}
\begin{itemize}
\raggedright
  \item \emph{Say it:} \emph{vajā māḍu}
  \item \emph{Say it:} \emph{vajā māḍu}, once more
  \item \emph{Say it:} \emph{taḍavāgi baruttidda naukaranannu kampani vajā māḍitu.} — The company sacked the employee who kept coming late.
  \item \emph{Recall:} say the Kannada for to gain, then the Kannada for to appoint, then say \emph{vajā māḍu} again
\end{itemize}

\begin{tcolorbox}[breakable,colback=teal!4,colframe=teal!35!black,title={Before you move on}]
What does \kn{ವಜಾ} \kn{ಮಾಡು} mean? (\textquotedblleft{}To dismiss; to sack\textquotedblright{}.) Say it once more.
\end{tcolorbox}
\section[\texorpdfstring{\kn{ನಿರುದ್ಯೋಗ} (nirudyōga)}{ನಿರುದ್ಯೋಗ (nirudyōga)}]{\kn{ನಿರುದ್ಯೋಗ} (nirudyōga) — unemployment}

\label{lesson:KA-C518-nirudyoga}



\begin{quote}
Before the new one: say the Kannada for to appoint, then the Kannada for to dismiss.
\end{quote}

\subsection*{You'll want to know: \kn{ನಿರುದ್ಯೋಗ}}
\textbf{\kn{ನಿರುದ್ಯೋಗ}} — \emph{nirudyōga} — \textquotedblleft{}unemployment\textquotedblright{}.

\begin{quote}
\kn{ಹಳ್ಳಿಗಳಲ್ಲಿ} \kn{ನಿರುದ್ಯೋಗ} \kn{ದೊಡ್ಡ} \kn{ಸಮಸ್ಯೆ}.
\end{quote}

\begin{quote}
\emph{haḷḷigaḷalli nirudyōga doḍḍa samasye.}
\end{quote}

\begin{quote}
Unemployment is a big problem in the villages.
\end{quote}

\kn{ನಿರ್}, without + \kn{ಉದ್ಯೋಗ}, a job. \kn{ನಿರುದ್ಯೋಗಿ}, an unemployed person.

\begin{cousinweb}
\textbf{\kn{ನಿರುದ್ಯೋಗ}} (\emph{nirudyōga}) is a learned compound built from Sanskrit \emph{nir-}, \textquotedblleft{}without\textquotedblright{}, and \emph{udyoga}, \textquotedblleft{}undertaking, exertion\textquotedblright{} (Monier-Williams).
\end{cousinweb}

\begin{grammarlens}[title={What you've built}]
One more word for describing what you lived and felt.
\end{grammarlens}

\subsection*{Your turn}
\begin{itemize}
\raggedright
  \item \emph{Say it:} \emph{nirudyōga}
  \item \emph{Say it:} \emph{nirudyōga}, once more
  \item \emph{Say it:} \emph{haḷḷigaḷalli nirudyōga doḍḍa samasye.} — Unemployment is a big problem in the villages.
  \item \emph{Recall:} say the Kannada for to appoint, then the Kannada for to dismiss, then say \emph{nirudyōga} again
\end{itemize}

\begin{tcolorbox}[breakable,colback=teal!4,colframe=teal!35!black,title={Before you move on}]
What does \kn{ನಿರುದ್ಯೋಗ} mean? (\textquotedblleft{}Unemployment\textquotedblright{}.) Say it once more.
\end{tcolorbox}
\section[\texorpdfstring{\kn{ನಿವೃತ್ತಿ} (nivṛtti)}{ನಿವೃತ್ತಿ (nivṛtti)}]{\kn{ನಿವೃತ್ತಿ} (nivṛtti) — retirement}

\label{lesson:KA-C518-nivrtti}



\begin{quote}
Before the new one: say the Kannada for to dismiss, then the Kannada for unemployment.
\end{quote}

\subsection*{You'll want to know: \kn{ನಿವೃತ್ತಿ}}
\textbf{\kn{ನಿವೃತ್ತಿ}} — \emph{nivṛtti} — \textquotedblleft{}retirement\textquotedblright{}.

\begin{quote}
\kn{ಅರವತ್ತನೇ} \kn{ವಯಸ್ಸಿನಲ್ಲಿ} \kn{ಅಪ್ಪ} \kn{ನಿವೃತ್ತಿ} \kn{ಪಡೆದರು}.
\end{quote}

\begin{quote}
\emph{aravattanē vayassinalli appa nivṛtti paḍedaru.}
\end{quote}

\begin{quote}
Father retired at the age of sixty.
\end{quote}

\kn{ನಿವೃತ್ತ}, retired.

\begin{cousinweb}
\textbf{\kn{ನಿವೃತ್ತಿ}} (\emph{nivṛtti}) is borrowed from Sanskrit \emph{nivṛtti}, \textquotedblleft{}returning; ceasing, abstaining\textquotedblright{} (Monier-Williams), from \emph{ni-}, \textquotedblleft{}back\textquotedblright{}, and \emph{vṛt}, \textquotedblleft{}to turn\textquotedblright{}.
\end{cousinweb}

\begin{grammarlens}[title={What you've built}]
One more word for describing what you lived and felt.
\end{grammarlens}

\subsection*{Your turn}
\begin{itemize}
\raggedright
  \item \emph{Say it:} \emph{nivṛtti}
  \item \emph{Say it:} \emph{nivṛtti}, once more
  \item \emph{Say it:} \emph{aravattanē vayassinalli appa nivṛtti paḍedaru.} — Father retired at the age of sixty.
  \item \emph{Recall:} say the Kannada for to dismiss, then the Kannada for unemployment, then say \emph{nivṛtti} again
\end{itemize}

\begin{tcolorbox}[breakable,colback=teal!4,colframe=teal!35!black,title={Before you move on}]
What does \kn{ನಿವೃತ್ತಿ} mean? (\textquotedblleft{}Retirement\textquotedblright{}.) Say it once more.
\end{tcolorbox}
\section[\texorpdfstring{\kn{ಗುತ್ತಿಗೆ} (guttige)}{ಗುತ್ತಿಗೆ (guttige)}]{\kn{ಗುತ್ತಿಗೆ} (guttige) — a contract; a lease}

\label{lesson:KA-C518-guttige}



\begin{quote}
Before the new one: say the Kannada for unemployment, then the Kannada for retirement.
\end{quote}

\subsection*{You'll want to know: \kn{ಗುತ್ತಿಗೆ}}
\textbf{\kn{ಗುತ್ತಿಗೆ}} — \emph{guttige} — \textquotedblleft{}a contract; a lease\textquotedblright{}.

\begin{quote}
\kn{ರಸ್ತೆ} \kn{ಕೆಲಸದ} \kn{ಗುತ್ತಿಗೆ} \kn{ಒಂದು} \kn{ಕಂಪನಿಗೆ} \kn{ಸಿಕ್ಕಿತು}.
\end{quote}

\begin{quote}
\emph{raste kelasada guttige ondu kampanige sikkitu.}
\end{quote}

\begin{quote}
The road-building contract went to a company.
\end{quote}

\kn{ಗುತ್ತಿಗೆದಾರ}, a contractor.

\begin{cousinweb}
\textbf{\kn{ಗುತ್ತಿಗೆ}} (\emph{guttige}) is a Kannada word for a contract or lease; its origin is not settled.
\end{cousinweb}

\begin{grammarlens}[title={What you've built}]
One more word for describing what you lived and felt.
\end{grammarlens}

\subsection*{Your turn}
\begin{itemize}
\raggedright
  \item \emph{Say it:} \emph{guttige}
  \item \emph{Say it:} \emph{guttige}, once more
  \item \emph{Say it:} \emph{raste kelasada guttige ondu kampanige sikkitu.} — The road-building contract went to a company.
  \item \emph{Recall:} say the Kannada for unemployment, then the Kannada for retirement, then say \emph{guttige} again
\end{itemize}

\begin{tcolorbox}[breakable,colback=teal!4,colframe=teal!35!black,title={Before you move on}]
What does \kn{ಗುತ್ತಿಗೆ} mean? (\textquotedblleft{}A contract; a lease\textquotedblright{}.) Say it once more.
\end{tcolorbox}
\section[(dialogue)]{Looking back: from Serious Illness to Hiring and Retiring}

\label{lesson:KA-R518-review-512-518}



\begin{quote}
Say the Kannada for blood pressure and a heart attack.
\end{quote}

\subsection*{Your turn}
Say these aloud:

\begin{itemize}
\raggedright
  \item blood pressure, a heart attack, bleeding, a stroke, to faint
  \item to recover, to feed, a heartbeat, nursing, a herbal decoction
  \item a diet, a vitamin, to waste away, to get tired, lifespan
  \item a blind man, deaf, mute, to groan, to droop
  \item a profession, a post, a qualification, an expert, expertise
  \item a promotion, progress, ambition, to achieve, to gain
  \item to appoint, to dismiss, unemployment, retirement, a contract
\end{itemize}

\begin{tcolorbox}[breakable,colback=teal!4,colframe=teal!35!black,title={Before you move on}]
Blood pressure? (\textbf{\kn{ರಕ್ತದೊತ್ತಡ}}, \emph{raktadottaḍa}.) A heart attack? (\textbf{\kn{ಹೃದಯಾಘಾತ}}, \emph{hṛdayāghāta}.) Bleeding? (\textbf{\kn{ರಕ್ತಸ್ರಾವ}}, \emph{raktasrāva}.) A stroke? (\textbf{\kn{ಪಾರ್ಶ್ವವಾಯು}}, \emph{pārśvavāyu}.) To faint? (\textbf{\kn{ಪ್ರಜ್ಞೆ} \kn{ತಪ್ಪು}}, \emph{prajñe tappu}.) To recover? (\textbf{\kn{ಚೇತರಿಸಿಕೊ}}, \emph{cētarisiko}.) To feed? (\textbf{\kn{ಉಣಿಸು}}, \emph{uṇisu}.) A heartbeat? (\textbf{\kn{ಬಡಿತ}}, \emph{baḍita}.) Nursing? (\textbf{\kn{ಆರೈಕೆ}}, \emph{āraike}.) A herbal decoction? (\textbf{\kn{ಕಷಾಯ}}, \emph{kaṣāya}.) A diet? (\textbf{\kn{ಪಥ್ಯ}}, \emph{pathya}.) A vitamin? (\textbf{\kn{ಜೀವಸತ್ವ}}, \emph{jīvasatva}.) To waste away? (\textbf{\kn{ಸೊರಗು}}, \emph{soragu}.) To get tired? (\textbf{\kn{ದಣಿ}}, \emph{daṇi}.) Lifespan? (\textbf{\kn{ಆಯುಷ್ಯ}}, \emph{āyuṣya}.) A blind man? (\textbf{\kn{ಕುರುಡ}}, \emph{kuruḍa}.) Deaf? (\textbf{\kn{ಕಿವುಡ}}, \emph{kivuḍa}.) Mute? (\textbf{\kn{ಮೂಕ}}, \emph{mūka}.) To groan? (\textbf{\kn{ನರಳು}}, \emph{naraḷu}.) To droop? (\textbf{\kn{ನಲುಗು}}, \emph{nalugu}.) A profession? (\textbf{\kn{ವೃತ್ತಿ}}, \emph{vṛtti}.) A post? (\textbf{\kn{ಹುದ್ದೆ}}, \emph{hudde}.) A qualification? (\textbf{\kn{ಅರ್ಹತೆ}}, \emph{arhate}.) An expert? (\textbf{\kn{ತಜ್ಞ}}, \emph{tajña}.) Expertise? (\textbf{\kn{ಪರಿಣತಿ}}, \emph{pariṇati}.) A promotion? (\textbf{\kn{ಬಡ್ತಿ}}, \emph{baḍti}.) Progress? (\textbf{\kn{ಪ್ರಗತಿ}}, \emph{pragati}.) Ambition? (\textbf{\kn{ಮಹತ್ವಾಕಾಂಕ್ಷೆ}}, \emph{mahatvākāṅkṣe}.) To achieve? (\textbf{\kn{ಸಾಧಿಸು}}, \emph{sādhisu}.) To gain? (\textbf{\kn{ಗಳಿಸು}}, \emph{gaḷisu}.) To appoint? (\textbf{\kn{ನೇಮಿಸು}}, \emph{nēmisu}.) To dismiss? (\textbf{\kn{ವಜಾ} \kn{ಮಾಡು}}, \emph{vajā māḍu}.) Unemployment? (\textbf{\kn{ನಿರುದ್ಯೋಗ}}, \emph{nirudyōga}.) Retirement? (\textbf{\kn{ನಿವೃತ್ತಿ}}, \emph{nivṛtti}.) A contract? (\textbf{\kn{ಗುತ್ತಿಗೆ}}, \emph{guttige}.)
\end{tcolorbox}
